\hypertarget{python-2.x-low-level-file-io-exceptions-unwinding-blocks}{%
\chapter{Python 2.x: Low-Level File I/O \& Exceptions Unwinding
Blocks}\label{python-2.x-low-level-file-io-exceptions-unwinding-blocks}}

\hypertarget{low-level-file-management-and-standard-streams-wrapping}{%
\subsection{6.1 Low-Level File Management and Standard Streams
wrapping}\label{low-level-file-management-and-standard-streams-wrapping}}

Under Python 2.x, the built-in \passthrough{\lstinline!file!} type and
\passthrough{\lstinline!open()!} function were thin wrappers around
standard C stdio library streams (\passthrough{\lstinline!FILE *!}).

\hypertarget{the-cpython-2.x-pyfileobject-layout}{%
\subsubsection{\texorpdfstring{1. The CPython 2.x \texttt{PyFileObject}
Layout}{1. The CPython 2.x PyFileObject Layout}}\label{the-cpython-2.x-pyfileobject-layout}}

Let's inspect the \passthrough{\lstinline!PyFileObject!} C structure
defined in CPython 2.x's \passthrough{\lstinline!Include/fileobject.h!}:

\begin{lstlisting}[language=C]
typedef struct {
    PyObject_HEAD
    FILE *f_fp;                 /* Pointer to standard C stdio stream */
    PyObject *f_name;           /* File name string object */
    PyObject *f_mode;           /* File mode string object */
    int (*f_close)(FILE *);     /* Close wrapper function */
    int f_softspace;            /* Print statement space tracking state */
    int f_binary;               /* Binary mode flag indicator */
    PyObject *f_encoding;       /* File encoding string */
    PyObject *f_errors;         /* Error handling mode */
    PyObject *f_newlines;       /* Decoded newlines object */
} PyFileObject;
\end{lstlisting}

\hypertarget{the-f_softspace-printing-flag}{%
\subsubsection{\texorpdfstring{2. The \texttt{f\_softspace} Printing
Flag}{2. The f\_softspace Printing Flag}}\label{the-f_softspace-printing-flag}}

In Python 2.x, the print statement was keyword-based rather than a
function. The runtime tracked layout separation using the
\passthrough{\lstinline!f\_softspace!} slot inside the standard output
stream: * When a print statement completed printing an item, it
evaluated the next token. If a trailing comma was present (e.g.,
\passthrough{\lstinline!print x,!}), the VM set
\passthrough{\lstinline!f\_softspace = 1!} and omitted the newline
character. * During the subsequent print operation, CPython checked
\passthrough{\lstinline!f\_softspace!}. If it was
\passthrough{\lstinline!1!}, it wrote a space to the file stream before
printing the next element and reset
\passthrough{\lstinline!f\_softspace = 0!}.

\hypertarget{redundant-locking-bottlenecks-under-the-gil}{%
\subsubsection{3. Redundant Locking Bottlenecks under the
GIL}\label{redundant-locking-bottlenecks-under-the-gil}}

Because Python 2.x wrapped standard C \passthrough{\lstinline!FILE *!}
streams, I/O operations were subject to the underlying C library's
internal buffering and synchronization locking. * Standard C libraries
acquire internal mutex locks per read/write call to ensure
multi-threaded safety. * In CPython, since the Global Interpreter Lock
(GIL) already guarantees single-threaded interpreter execution, these
nested C stdio locks were redundant, introducing unnecessary CPU
context-switching overhead and locking bottlenecks during high-frequency
concurrent file operations.

\hypertarget{the-layered-io-system-pep-3116}{%
\subsubsection{4. The Layered I/O System (PEP
3116)}\label{the-layered-io-system-pep-3116}}

To eliminate stdio bottlenecks, Python 3.0 replaced the old stream-based
system with a layered I/O architecture backported to Python 2.6+. This
system interacts directly with OS file descriptors via system calls
(e.g., \passthrough{\lstinline!read(2)!},
\passthrough{\lstinline!write(2)!}), bypassing C's stdio library: 1.
\textbf{Raw Layer (\passthrough{\lstinline!RawIOBase!} /
\passthrough{\lstinline!FileIO!})}: A thin wrapper over the POSIX OS
file descriptor, executing raw system-level reads and writes. 2.
\textbf{Buffered Layer (\passthrough{\lstinline!BufferedIOBase!} /
\passthrough{\lstinline!BufferedReader!} /
\passthrough{\lstinline!BufferedWriter!})}: Maintains an internal memory
block (typically 8KB) to minimize context switches between user space
and kernel space. 3. \textbf{Text Layer
(\passthrough{\lstinline!TextIOBase!} /
\passthrough{\lstinline!TextIOWrapper!})}: Handles encoding/decoding and
system-specific universal newline translations.

\begin{center}\rule{0.5\linewidth}{0.5pt}\end{center}

\hypertarget{exception-mechanics-and-the-thread-state-slots}{%
\subsection{6.2 Exception Mechanics and the Thread State
slots}\label{exception-mechanics-and-the-thread-state-slots}}

CPython manages exceptions at the thread level, storing active and
caught exception states inside the execution thread's state structure.

\hypertarget{exception-slots-in-pythreadstate}{%
\subsubsection{\texorpdfstring{1. Exception Slots in
\texttt{PyThreadState}}{1. Exception Slots in PyThreadState}}\label{exception-slots-in-pythreadstate}}

CPython maintains exception metadata on a per-thread basis within the
\passthrough{\lstinline!PyThreadState!} struct
(\passthrough{\lstinline!Include/pystate.h!}):

\begin{lstlisting}[language=C]
typedef struct _ts {
    struct _ts *next;
    PyInterpreterState *interp;
    struct _frame *frame;       /* Active execution frame */
    
    /* Exception currently propagating */
    PyObject *curexc_type;
    PyObject *curexc_value;
    PyObject *curexc_traceback;

    /* Exception currently caught and handled */
    PyObject *exc_type;
    PyObject *exc_value;
    PyObject *exc_traceback;
} PyThreadState;
\end{lstlisting}

\begin{itemize}
\tightlist
\item
  \textbf{\passthrough{\lstinline!curexc\_*!} (Current Exception)}: The
  active exception currently propagating. When an instruction fails
  (e.g., division by zero), CPython sets these pointers and returns
  \passthrough{\lstinline!NULL!} up the C execution chain.
\item
  \textbf{\passthrough{\lstinline!exc\_*!} (Caught Exception)}: The
  exception currently being handled inside an active
  \passthrough{\lstinline!except!} block. This keeps the exception
  accessible via \passthrough{\lstinline!sys.exc\_info()!} during
  nesting, while preventing active propagation from overwriting the
  caught state.
\end{itemize}

\hypertarget{the-string-exception-era}{%
\subsubsection{2. The String Exception
Era}\label{the-string-exception-era}}

In early Python versions (pre-2.6), exceptions were not required to be
class instances; you could raise plain string literals:

\begin{lstlisting}[language=Python]
# Pre-Python 2.6 Behavior
raise "DatabaseConnectionFailed"
\end{lstlisting}

During lookup, the VM evaluated string exceptions by identity
(\passthrough{\lstinline!is!} check) rather than class inheritance. This
made hierarchy categorization and error composition difficult. Python
2.6 deprecated and removed string exceptions, requiring all exceptions
to inherit from the built-in base type
\passthrough{\lstinline!BaseException!}.

\begin{center}\rule{0.5\linewidth}{0.5pt}\end{center}

\hypertarget{the-try-except-finally-block-stack-and-stack-restoration}{%
\subsection{6.3 The Try-Except-Finally Block Stack and Stack
Restoration}\label{the-try-except-finally-block-stack-and-stack-restoration}}

CPython handles structured control flow (loops, exception blocks) using
a static \textbf{block stack} located inside each execution frame
(\passthrough{\lstinline!PyFrameObject!}).

\hypertarget{the-pytryblock-structure}{%
\subsubsection{\texorpdfstring{1. The \texttt{PyTryBlock}
Structure}{1. The PyTryBlock Structure}}\label{the-pytryblock-structure}}

The block stack consists of up to 20
(\passthrough{\lstinline!CO\_MAXBLOCKS!})
\passthrough{\lstinline!PyTryBlock!} structs:

\begin{lstlisting}[language=C]
typedef struct {
    int b_type;                 /* Block type: SETUP_LOOP, SETUP_EXCEPT, SETUP_FINALLY */
    int b_handler;              /* Target instruction offset of the handler */
    int b_level;                /* Stack pointer depth at block entry */
} PyTryBlock;
\end{lstlisting}

\hypertarget{vm-exception-unwinding-flow}{%
\subsubsection{2. VM Exception Unwinding
Flow}\label{vm-exception-unwinding-flow}}

When an instruction returns \passthrough{\lstinline!NULL!}, the VM loop
(\passthrough{\lstinline!\_PyEval\_EvalFrameDefault!}) enters unwinding
mode:

\begin{lstlisting}
[Exception Raised: C API returns NULL]
                  |
                  v
[Is frame->f_blockstack empty?]
        |             |
       Yes            No
        |             |
        v             v
[Pop Frame]     [Pop Try Block]
[Traceback]           |
        |       [Is block SETUP_EXCEPT or SETUP_FINALLY?]
        v             |                  |
[Repeat in            No                Yes
 caller]              |                  |
                      v                  v
                [Continue loop]   1. Restore stack pointer: stack_pointer = frame->b_level
                                  2. Push traceback, value, type onto stack
                                  3. Set PC: frame->f_lasti = block->b_handler
                                  4. Resume bytecode execution in handler
\end{lstlisting}

\begin{enumerate}
\def\labelenumi{\arabic{enumi}.}
\tightlist
\item
  \textbf{Stack Cleanup via \passthrough{\lstinline!b\_level!}}: The VM
  pops the current try block. It immediately restores the frame's
  evaluation stack pointer back to the offset saved in
  \passthrough{\lstinline!b\_level!}. This discards any unused variables
  or intermediate states created inside the
  \passthrough{\lstinline!try!} block, preventing memory leaks.
\item
  \textbf{Context Setup}: The VM pushes the exception's
  \passthrough{\lstinline!traceback!}, \passthrough{\lstinline!value!},
  and \passthrough{\lstinline!type!} (in Python 2.x) onto the evaluation
  stack.
\item
  \textbf{Frame Traversal fallback}: If the frame's block stack is
  exhausted without finding a handler, CPython pops the frame,
  instantiates a \passthrough{\lstinline!PyTracebackObject!} mapping the
  exception to the current line number, and continues unwinding in the
  caller's frame (\passthrough{\lstinline!frame->f\_back!}).
\end{enumerate}

\begin{center}\rule{0.5\linewidth}{0.5pt}\end{center}
